\chapter{Apresentação}

O presente Projeto Pedagógico do Curso (PPC) de \nomecursonovo da Universidade Federal de Uberlândia foi elaborado em 2025 com o propósito de consolidar uma formação acadêmica inovadora, crítica e socialmente engajada, alinhada às Diretrizes Curriculares Nacionais (DCNs) da área de Computação~\cite{MEC:CNE:CES:5:2016} e às orientações institucionais da UFU~\cite{UFU:CONGRAD:15:2016,UFU:GuiaPPC:2021}. Esta atualização marca um novo ciclo de planejamento pedagógico, em consonância com as transformações contemporâneas que atravessam a ciência, a tecnologia, a educação e o mundo do trabalho.

Desde a última reformulação em 2018, muitos avanços e desafios se impuseram: o impacto da pandemia de COVID-19 na organização do ensino superior, a consolidação da curricularização da extensão como diretriz nacional (meta 12.7 do PNE)~\cite{mec2014pne,UFU:CONSUN:25:2019,UFU:CONGRAD:13:2019,UFU:CONSEX:5:2020}, a ampliação das possibilidades de ensino híbrido~\cite{Decreto:12456:2025}, e a crescente demanda por internacionalização, diversidade e inovação nos currículos~\cite{UFU:CONSUN:31:2022}. Em resposta a esse cenário, este novo PPC foi construído de forma coletiva pelo Núcleo Docente Estruturante (NDE), pelo Colegiado do Curso e por representantes da comunidade acadêmica, com a convicção de que uma formação em Engenharia de Computação deve ser dinâmica, plural e orientada por competências.

A revisão curricular de 2025 reafirma os princípios fundadores do curso — como a centralidade da programação contínua, a exposição a múltiplos paradigmas computacionais, e a valorização da pesquisa docente — ao mesmo tempo em que incorpora diretrizes contemporâneas como a educação por competências (CBL), a flexibilização do percurso formativo, a inclusão efetiva de práticas extensionistas, a valorização da diversidade e o estímulo à formação integral dos discentes. Além disso, foram atualizadas as referências internacionais utilizadas como parâmetro para a construção do currículo, como o \textit{Computer Curricula 2020} (CC2020 da ACM/IEEE)~\cite{ACM:CC2020} e os princípios\footnote{Relacionados aos Objetivos de Desenvolvimento Sustentável da ONU, disponível em: \url{https://brasil.un.org/pt-br/sdgs}} do \textit{Engineering Education for Sustainable Development} (EESD)~\cite{UNESCO:Eng4SusDev:2021}.

Este PPC é, assim, expressão de um compromisso institucional com a excelência acadêmica, a inovação formativa, o impacto social e o protagonismo dos estudantes. Mais do que uma exigência normativa, ele representa a carta de intenções de um curso que se projeta para o futuro, sem perder de vista suas raízes e responsabilidades com o presente.